\section{Related Work}

\textbf{Long Context LLMs} Extrapolation methods for RoPE-based LLMs~\cite{su2024roformer}, such as NTK~\cite{blocntkaware}, PI~\cite{chen2023extending}, YaRN~\cite{peng2023yarn} and DCA~\cite{an2024training}, modify the base frequency, position index and other components of positional embeddings, enabling the model to capture long-range semantic dependencies. On the other hand, Linear attention mechanisms~\cite{child2019generating, katharopoulos2020transformers}, Recurrent Neural Networks (RNNs) and State Space Models (SSMs)~\cite{gu2021efficiently, gu2023mamba, peng2023rwkv, de2024griffin, feng2024attention} leverage different architectures to achieve $O(N)$ computation complexity, aiming to process extremely long context. Sparse attention ~\cite{beltagy2020longformer, zhao2019explicit, xiao2023efficient} shifts the attention mask matrix to apply patterns such as sliding windows, thereby eliminating irrelevant attention calculations. However, these patterns are typically based on predefined heuristics. The possibility of dynamic sparse attention~\cite{yuan2025native,lu2025moba} has been explored recently. The Long Short-Term Memory (LSTM) mechanism ~\cite{hochreiter1997long} achieved significant success in early NLP tasks, while Neural Turing Machines ~\cite{graves2014neural} and Memory Networks~\cite{weston2014memory} demonstrated how to equip neural networks with memory. Existing memory mechanisms integrated to Transformer models are typically realized by adding external memory modules~\cite{martins2021infty,wu2020memformer, behrouz2024titans, bulatov2023scaling} or external database~\cite{zhong2024memorybank, lu2023memochat, modarressi2023ret}. In contrast, we use reinforcement learning (RL) to enable LLM itself the ability to memorize.


\textbf{Reinforcement Learning for LLMs} In recent RL studies, the reward signals have gradually shifted from human preferences ~\cite{NEURIPS2022_b1efde53} or reward models distilled from them~\cite{bai2022constitutional} to rule-based feedback, which has demonstrated great potential in enhancing model reasoning capabilities~\cite{o1, guo2025deepseek, qwq, gemini-thinking, k1.5}. Key contributions include PPO ~\cite{schulman2017proximal} based on GAE~\cite{schulman2018highdimensionalcontinuouscontrolusing},  the Actor-Critic framework, as well as GRPO~\cite{deepseekmath} that utilizes Group Normalization. Algorithmic enhancements~\cite{hu2025reinforce++, yu2025dapo, liu2025understanding} have mostly focused on improving training sustainability and sample efficiency of these algorithms. To further release the potential of RL, recent works such as Search-R1~\cite{jin2025search}, Agent-R1~\cite{Agent-R1} and RAGEN~\cite{ragen} have explored the training of tool-using agents based on multi-turn chat. However, these multi-turn chats are constructed by alternately concatenating tool responses and model replies, with the ultimate optimization goal being a single conversation with tool masking. GiGPO~\cite{feng2025group} further investigates the use of multiple independent contexts in agent training with environment feedback with sliding window trajectories. However, these approaches are limited to optimizing interleaved trajectories of observation and generation, making them difficult to apply to more general agent workflows.

